% ==========================================================
% 00_config.tex
% Arquivo central de configuração do template LaTeX
% Organiza pacotes, metadados do documento e layout da página.
% ==========================================================

% ==========================================================
% METADADOS E VARIÁVEIS DO DOCUMENTO
% ==========================================================
\title{Projeto FEMM -- Máquinas de Corrente Alternada}                         % Título principal do documento.
\author{Luis Henrique Ribeiro Maciel}       % Nome do autor.
\date{Maio 2026}                           % Data exibida na capa.

\def\uname{Universidade Federal de Minas Gerais} % Nome da universidade.
\def\depart{Escola de Engenharia}                % Unidade ou departamento.
\def\materia{Conversão de Energia}               % Disciplina.
\def\turma{Turma TA}                              % Turma.

% Salva título, autor e data em comandos reutilizáveis.
\makeatletter
\let\thetitle\@title
\let\theauthor\@author
\let\thedate\@date
\makeatother

% ----------------------------------------------------------
% IDIOMA, ENCODIFICAÇÃO E FORMATAÇÃO BÁSICA DO TEXTO
% ----------------------------------------------------------
\usepackage[utf8]{inputenc}   % Permite escrever caracteres acentuados diretamente no arquivo.
\usepackage[brazil]{babel}    % Ajusta hifenização, nomes automáticos e convenções para português do Brasil.
\usepackage{indentfirst}      % Indenta o primeiro parágrafo de cada seção.
\usepackage{parskip}          % Adiciona espaçamento entre parágrafos para melhorar a legibilidade.

% ----------------------------------------------------------
% MATEMÁTICA, SÍMBOLOS E NOTAÇÃO TÉCNICA
% ----------------------------------------------------------
\usepackage{amsmath}          % Ambientes e comandos avançados para fórmulas matemáticas.
\usepackage{amssymb}          % Conjunto ampliado de símbolos matemáticos.
\usepackage{steinmetz}        % Facilita a escrita de notação fasorial e ângulos, útil em elétrica.
\usepackage{pifont}           % Disponibiliza símbolos extras, como marcas, setas e ícones simples.

% ----------------------------------------------------------
% FIGURAS, LEGENDAS E ELEMENTOS FLUTUANTES
% ----------------------------------------------------------
\usepackage{graphicx}         % Permite inserir imagens no documento.
\usepackage{caption}          % Personaliza legendas de figuras e tabelas.
\usepackage{subcaption}       % Permite criar subfiguras e sublegendas.
\usepackage{float}            % Oferece maior controle sobre a posição de figuras e tabelas.
\usepackage{tikz}             % Ferramenta poderosa para criar gráficos vetoriais diretamente no LaTeX, útil para diagramas de rede.
\usepackage{enumitem}         % Personaliza listas enumeradas e não enumeradas, como itemize e enumerate.
\usetikzlibrary{shapes.geometric, arrows.meta, positioning}

% ----------------------------------------------------------
% TABELAS E ESTRUTURA TABULAR
% ----------------------------------------------------------
\usepackage{booktabs}       % Melhora a qualidade visual de tabelas com regras horizontais adequadas.
\usepackage{multirow}       % Permite mesclar várias linhas em uma mesma célula de tabela.
\usepackage{tabularx}       % Oferece colunas de largura automática, útil para tabelas que precisam se ajustar ao texto.
\usepackage{array}          % Oferece mais opções para formatação de colunas em tabelas.
\usepackage{longtable}      % Permite criar tabelas que se estendem por várias páginas, útil para dados extensos.

% ----------------------------------------------------------
% AUTORIA E INFORMAÇÕES DE CAPA
% ----------------------------------------------------------
\usepackage{authblk}          % Facilita a organização de autores e afiliações em documentos acadêmicos.

% ----------------------------------------------------------
% CÓDIGO-FONTE E DESTAQUE DE SINTAXE
% ----------------------------------------------------------
\usepackage[dvipsnames]{xcolor}             % Define cores adicionais, úteis em links, código e elementos visuais.
\usepackage{minted}                         % Formatação de código com destaque de sintaxe via Pygments.
\usepackage{matlab-prettifier}              % Melhora a exibição de trechos de código MATLAB.
\usepackage{tcolorbox}                      % Criação de caixas de texto coloridas, úteis para destacar definições, exemplos ou código.
\tcbuselibrary{minted, skins, breakable}    % Habilita bibliotecas do tcolorbox para integração com minted e opções avançadas de formatação.

% ----------------------------------------------------------
% HIPERLINKS E NAVEGAÇÃO NO PDF
% ----------------------------------------------------------
\usepackage[colorlinks=true, allcolors=blue]{hyperref} % Cria links clicáveis no PDF, como sumário e referências.

% ----------------------------------------------------------
% CABEÇALHO, RODAPÉ E DIMENSÕES DA PÁGINA
% ----------------------------------------------------------

% Margens exigidas: 2,0 cm em todos os lados
\usepackage[
    a4paper,
    left=2cm,
    right=2cm,
    top=2cm,
    bottom=2cm,
    includehead
]{geometry}

% Espaçamento entre linhas exatamente 15 pt
\AtBeginDocument{\fontsize{12pt}{15pt}\selectfont}

% Espaço entre parágrafos exatamente 3 pt
\setlength{\parskip}{3pt}

% Recuo de parágrafo
\setlength{\parindent}{1.25cm}

% Cabeçalho e numeração no canto superior direito
\usepackage{fancyhdr}
\cfoot{\thepage} % Numeração centralizada no rodapé

% ----------------------------------------------------------
% ESTILO DA BIBLIOGRAFIA
% ----------------------------------------------------------
\bibliographystyle{plain}     % Define o estilo bibliográfico padrão simples.
